%!TeX root=../tese.tex
%("dica" para o editor de texto: este arquivo é parte de um documento maior)
% para saber mais: https://tex.stackexchange.com/q/78101

% As palavras-chave são obrigatórias, em português e em inglês, e devem ser
% definidas antes do resumo/abstract.
\palavraschave{evolução semântica, bancos de dados distribuídos, particionamento, replicação, MongoDB}

\keywords{semantic evolution, distributed databases, sharding, replication, MongoDB}

% O resumo é obrigatório, em português e inglês, em parágrafo único, com no
% máximo 500 palavras, apresentando objetivos, métodos, resultados e conclusões.
\resumo{
O MellowDB é um \textit{middleware} que gerencia a evolução semântica dos dados, as mudanças de significado ao longo do tempo, sobre um nó único do MongoDB, com duas estratégias de resolução, a reescrita de consultas (\textit{rewrite}) e o pré-processamento com versionamento (\textit{preprocess}). Este trabalho leva o MellowDB a ambientes distribuídos, a replicação e o particionamento, e mede o custo dessa transição. Implementou-se a replicação com a cadeia de versões autoritativa no primário, sob transação multidocumento, e as leituras distribuídas aos secundários, e o particionamento com a chave de \textit{shard}, a materialização das versões sobre uma coleção particionada e o tratamento de documentos sem o campo da chave. Uma campanha de \textit{benchmarks} em um \textit{cluster} na nuvem, sobre um corpus sintético de assimetria e cardinalidade controladas e um corpus real do DATASUS (transição da CID-9 para a CID-10), levou a duas conclusões. A primeira é que o compromisso entre as estratégias não tem vencedor fixo, pois o \textit{preprocess} predomina, mas a distribuição encurta a diferença e a profundidade da cadeia de evolução chega a invertê-la. A segunda é que, quando a chave de \textit{shard} recai sobre o campo que evolui, a operação realoca fisicamente os dados, e o desequilíbrio resultante pode ser irreparável ou impor um custo real de reequilíbrio, de modo que não é recomendado que a chave de \textit{shard} seja o campo que a evolução reescreve.
}

\abstract{
MellowDB is a \textit{middleware} that manages semantic evolution in data, the changes in meaning over time, on a single MongoDB node, with two resolution strategies, query rewriting (\textit{rewrite}) and preprocessing with versioning (\textit{preprocess}). This work brings MellowDB to distributed environments, replication and partitioning, and measures the cost of this transition. Replication was implemented with the version chain authoritative on the primary, under a multidocument transaction, and reads offloaded to the secondaries, and partitioning with the \textit{shard} key, the materialization of the versions over a partitioned collection, and the handling of documents without the key field. A \textit{benchmark} campaign on a cloud \textit{cluster}, over a synthetic corpus with controlled skew and cardinality and a real corpus from DATASUS (the transition from CID-9 to CID-10), led to two conclusions. The first is that the trade-off between the strategies has no fixed winner, since \textit{preprocess} prevails, but distribution narrows the gap and the depth of the evolution chain can even invert it. The second is that, when the \textit{shard} key falls on the evolving field, the operation physically relocates the data, and the resulting imbalance may be irreparable or impose a real rebalancing cost, so that it is not recommended that the \textit{shard} key be the field that evolution rewrites.
}
